\section{Method}
\label{sec:method}

\paragraph{Setup.}
All experiments used a rented NVIDIA RTX A4000 (16\,GB) with the NVIDIA 580-series
driver and vLLM 0.29.0, serving Qwen2.5-0.5B-Instruct. Sessions ran on several rented
virtual machines, each with its own physical GPU; Phase C reused Phase B's machine. The
client was \texttt{vllm bench serve} on the same host, issuing a closed-loop workload.
Up to a fixed number of requests (the concurrency, written C1 to C32) are kept in
flight, and a new request is sent as soon as one finishes. Each run sends 100 requests
(50 for the decode-heavy workload), so a run at C32 starts with a burst of 32 requests
and ends as the last ones drain. Prompts are synthetic, with exact input and output
lengths and a fixed seed.
% src: configs/*.toml (num_prompts 100; 50 for decode-heavy); reports/milestone1-*.md (driver 580.178.04); reports/milestone2-*.md (different physical GPU); reports/milestone3-*-phase-c.md (same VM as Phase B); reports/milestone5p2-*-part2.md (Provenance)

\paragraph{Controls.}
A first baseline (Milestone 1) compared C1 with C8, in forward and reverse order with
three repetitions and 15-second cooldowns, and found 6.65$\times$ higher output
throughput at C8. It ran with prefix caching on, an unpinned model revision, and the
model's own sampling defaults (temperature 0.7, top-p 0.8, top-k 20, repetition
penalty 1.1), any of which could inflate or mask the effect. A controlled repeat
(Milestone 2) turned prefix caching off, pinned the server's model and tokenizer revisions,
used greedy sampling, ignored the model's generation defaults
(\texttt{--generation-config vllm}), and warmed up at each measured concurrency. The
effect held (about 6.5$\times$), and these controls are fixed in every later run.
% src: reports/milestone1-20261001T205137Z.md (Result: 6.65x); reports/milestone2-20261003T034456Z.md (~6.5x); experiments/roadmap.md

\paragraph{Pre-registration.}
Each phase's protocol was merged into the public repository before its data were
collected. Milestones 1 and 2 stated expected signatures; from Milestone 3 on, each
protocol has predictions with thresholds and falsification rules (Phase B states them
as signatures for each workload regime). The merge
time is the evidence of pre-registration. Each Milestone 3 and 5 report scores every prediction as
confirmed, mixed, or not confirmed (with a few qualified verdicts, such as "not
testable"), and marks interpretations made after seeing the data as post hoc.
Section~\ref{sec:ledger} collects the 59 scored predictions of Milestones 3 and 5.

\paragraph{Three views of every run.}
From Milestone 3 on, each run is measured three ways.
\begin{itemize}
  \item \emph{Client:} time to first token (TTFT), the gaps between streamed tokens
    (inter-token latency, ITL; we report the median, p50, and the 99th percentile,
    p99), end-to-end latency (E2E), and output throughput.
  \item \emph{Server:} vLLM's Prometheus histograms for queue, prefill, and decode
    time per request, scraped before and after each run.
  \item \emph{Resources:} GPU utilization from \texttt{nvidia-smi} every 200\,ms,
    and CPU utilization every second.
\end{itemize}

\paragraph{A per-step scheduler log.}
vLLM's per-iteration token histogram credits a request's whole prompt to the step that
emits its first token, so it cannot show how many prompt tokens a step actually
processed. A manipulation check that relied on it failed as specified
(Section~\ref{sec:budget}). For Milestone 5 we added an opt-in vLLM plugin that wraps
the scheduler and writes one line per engine step, recording the scheduled tokens
(split into prefill and decode) and the scheduled requests. Part~2 added each prefill
chunk's context position and the time the step's results returned.
From these we compute a step's \emph{execution time}: the time from when it could
start (it was scheduled and the previous step's results had arrived) until its own
results arrived. Unlike the interval between scheduling calls, this measure is tied to
the step itself even when vLLM schedules the next step early. vLLM 0.29.0 can schedule
a step before the previous step's results return; we did not record which scheduling
mode was active.
% src: reports/milestone3-*-phase-d.md (iteration histogram); experiments/milestone5.md (Instrument; Part 2 addendum)

\paragraph{Validation and provenance.}
Every series is validated before analysis. The validator checks:
\begin{itemize}
  \item that runs completed with zero failures and exact token lengths;
  \item that the server's request count matches the prompts sent;
  \item that warm-ups ran and telemetry covers the series (sampling gaps of up to
    1\,s for the GPU and 3\,s for the CPU are allowed);
  \item for Milestone 5, that the step log covers every measured run.
\end{itemize}
Raw results are archived with a SHA-256 checksum, verified after transfer, before the
machine is deleted. Across all sessions the study comprises 18{,}900 measured requests.
% src: measured-request counts from each report's Validation table: M1 1,200; M2 1,200; B 1,800; C 4,800; D 5,400; M5 pilot 2,100; M5 Part 2 2,400
